% !TeX root = ../main.tex
\chapter{Einleitung}
\label{cha:einleitung}

\section{Motivation, Ziel und Abgrenzung der Arbeit}
\label{sec:motivation}
% Quelle: das Exposé der Aufgabenstellung. Ausgangsproblem, Ziel der Arbeit und
% was ausdrücklich NICHT Teil der Arbeit ist.
%
% Arbeit mit bekannten Anforderungen: Abschnitt heißt "Motivation und Ziel der
% Arbeit", danach folgt ein eigener Abschnitt "Anforderungen und Abgrenzung",
% und das Kapitel Anforderungsanalyse entfällt.

\section{Verwandte Arbeiten}
\label{sec:verwandte-arbeiten}
% Literaturrecherche (Bibliothekskatalog) und Internetrecherche. Suchmaschinen,
% Suchbegriffe und relevante Treffer dokumentieren. Eine Tabelle
% Suchmaschine | Suchbegriffe | relevante Treffer hat sich bewährt.

\section{Inhaltlicher Aufbau der Arbeit}
\label{sec:aufbau}
% Wenige Zeilen je Folgekapitel, mit \autoref auf die Kapitel.
